\chapter{Background and Related Work}
\label{sec:background}

\section{What Is Spatial Intelligence?}

Spatial intelligence is not a single capability but a family of related skills. Following the taxonomy proposed in~\cite{cai2025holisticevaluationmultimodalllms}, we decompose it into six fundamental dimensions.

\textbf{Metric Measurement (MM)} involves understanding physical scale such as estimating distances between objects, gauging the size of a room, or judging whether a sofa would fit through a doorway. \textbf{Spatial Relations (SR)} captures the ability to reason about relative positions of multiple objects within a 3D coordinate system. \textbf{Mental Reconstruction (MR)} requires inferring a detailed spatial representation of an object from limited 2D views. \textbf{Perspective-Taking (PT)} is perhaps the most demanding as it requires reasoning about how a scene appears from different viewpoints, establishing correspondences across views, inferring camera motion, and mentally simulating viewpoint shifts. \textbf{Deformation and Assembly (DA)} addresses scenarios where objects change shape or must be combined, for example understanding how a flat sheet folds into a box, tracing the steps of tying a knot, or figuring out how separate components fit together into a whole. Finally, \textbf{Comprehensive Reasoning (CR)} involves coordinating multiple spatial capabilities across multi-step reasoning chains.

This taxonomy proved essential for the work that followed. Without it, evaluation results would be a soup of numbers from disparate benchmarks; with it, we could precisely identify which spatial capabilities are strong, which are weak, and where targeted intervention is most needed.

\section{The Benchmark Landscape}

A growing ecosystem of benchmarks has emerged to test these capabilities, each covering a different slice of the space. VSI-Bench~\cite{yang2025thinkingspacemultimodallarge} evaluates video-based spatial reasoning (object counting, distance estimation, route planning) using egocentric indoor video. MMSI~\cite{yang2025mmsibenchbenchmarkmultiimagespatial} tests multi-image spatial understanding, focusing on positional relationships between cameras, objects, and regions. MindCube~\cite{yin2025spatialmentalmodelinglimited} probes whether models can infer what lies beyond the visible frame, such as filling in unobserved regions, adopting novel viewpoints, and predicting the outcome of hypothetical movements given only a few reference images. ViewSpatial~\cite{li2025viewspatialbenchevaluatingmultiperspectivespatial} probes view-dependent spatial reasoning. SITE~\cite{wang2025sitespatialintelligencethorough} provides a multi-choice VQA benchmark spanning single-image, multi-image, and video modalities, testing spatial visualization, orientation, and dynamic scene reasoning at both figural and environmental scales. OmniSpatial~\cite{jia2026omnispatialcomprehensivespatialreasoning} spans dynamic reasoning, spatial interaction, complex logic, and perspective-taking.

The fragmentation is the problem. Each benchmark uses its own evaluation protocol, prompt format, and metric. A model might score well on VSI-Bench's object counting but poorly on its route planning; it might handle MMSI's camera-camera correspondences but fail on object-object ones. Without a unified framework, comparing models across the full breadth of spatial capabilities is impractical.

\section{Approaches to Improving Spatial Reasoning}

Efforts to close the spatial intelligence gap follow two broad strategies. Architecture-based approaches incorporate 3D-aware encoders. Spatial-MLLM~\cite{wu2025spatialmllmboostingmllmcapabilities} pairs a standard 2D encoder with a visual geometry foundation model to extract 3D structural features from 2D inputs, VLM-3R~\cite{fan2025vlm3rvisionlanguagemodelsaugmented} derives implicit 3D tokens from monocular video via a geometry encoder, 3DThinker~\cite{chen2026think3dgeometricimagination} teaches VLMs to generate 3D latent representations during reasoning by aligning them with a 3D foundation model, then refining through outcome-based reinforcement learning. Data-centric approaches instead curate spatial training data and apply supervised fine-tuning (SFT) or reinforcement learning (RL). SpatialVLM~\cite{chen2024spatialvlmendowingvisionlanguagemodels} pioneered this direction with 2 billion synthetic spatial QA samples. Subsequent work includes SpatialLadder (26K samples with progressive training)~\cite{li2025spatialladderprogressivetrainingspatial}, VST (4.1M SFT + 135K RL samples)~\cite{yang2025visualspatialtuning}, and Cambrian-S (590K spatial video samples with four-stage training)~\cite{yang2025cambriansspatialsupersensingvideo}.

The data-centric approach is appealing because it does not require modifying model architectures, meaning any foundation model can potentially benefit. But the question of \textit{which} data, \textit{how much}, and \textit{organized how} remained open.

\section{From Static Benchmarks to Embodied Agents}

Static benchmarks, however comprehensive, test spatial intelligence in a fundamentally passive mode where the model receives images or video frames and answers questions. But spatial intelligence exists for action. Humans develop spatial reasoning not as an academic exercise but because it enables us to navigate rooms, manipulate objects, and coordinate with others in shared physical spaces.

Embodied AI benchmarks close this gap by placing agents in simulated 3D environments where they must perceive, reason, plan, and act. EmbodiedBench~\cite{yang2025embodiedbenchcomprehensivebenchmarkingmultimodal} provides a multi-domain evaluation spanning household tasks, navigation, habitat exploration, and robotic manipulation. HAZARD~\cite{zhou2024hazardchallengeembodieddecision} places agents in emergency scenarios (fire, flood, wind) where the goal is not explicitly specified as a point-to-point instruction but must be inferred from context, testing whether agents can reason about implicit objectives like prioritising and evacuating valuable objects from danger. VLN-CE~\cite{krantz2020navgraphvisionandlanguagenavigationcontinuous} requires translating natural language navigation instructions into continuous movement through realistic indoor environments. ManipulaTHOR~\cite{ehsani2021manipulathorframeworkvisualobject} tests arm-based point navigation requiring precise spatial positioning.

These benchmarks are expensive to run (each requiring its own simulator, conda environment, and GPU resources) and have historically been evaluated in isolation. No unified framework existed for running them all under a common protocol.
